\section{The generalized lambda distribution}\label{sec:background}

\subsection{Two parameterizations}

The GLD is defined by its quantile function. In the Ramberg-Schmeiser (RS)
form \parencite{RambergSchmeiser1974},
\begin{equation}
    Q_{\mathrm{RS}}(u) = \lambda_1 + \frac{u^{\lambda_3} - (1-u)^{\lambda_4}}{\lambda_2},
    \qquad 0 < u < 1,
    \label{eq:qrs}
\end{equation}
while the FKML form of \textcite{FreimerEtAl1988} writes
\begin{equation}
    Q_{\mathrm{FKML}}(u) = \lambda_1 +
    \frac{\varphi(u; \lambda_3) - \varphi(1-u; \lambda_4)}{\lambda_2},
    \qquad
    \varphi(u; \lambda) =
    \begin{cases}
        \dfrac{u^{\lambda} - 1}{\lambda}, & \lambda \neq 0, \\[8pt]
        \ln u,                            & \lambda = 0.
    \end{cases}
    \label{eq:qfkml}
\end{equation}
In both, $\lambda_1$ is a location parameter, $\lambda_2$ an inverse scale
parameter, and $\lambda_3, \lambda_4$ control the left and the right tail
separately. In the FKML form $\lambda_2 > 0$, whereas in the RS form the sign of
$\lambda_2$ is tied to the signs of the shape parameters, as
Section~\ref{subsec:validity} details. The two families coincide at
$\lambda_3 = \lambda_4 = 1$, where both reduce to a uniform distribution. The
FKML form is the reparameterization that removes the restrictions the RS form
places on the joint signs of the shape parameters, at the cost of requiring the
logarithmic limit in Equation~\eqref{eq:qfkml}.

Because $Q$ is available in closed form and its inverse is not, the natural
operations are sampling and quantile evaluation: a variate is drawn as
$Q(U)$ with $U \sim \mathrm{Uniform}(0,1)$, at the cost of one uniform draw. The
density is available only in the quantile-density form
\begin{equation}
    f(Q(u)) = \frac{1}{Q'(u)},
    \label{eq:density}
\end{equation}
so evaluating $f$ or $F$ at a given $x$ requires inverting $Q$ numerically.

\subsection{Support and moments}

The support $S = [Q(0^{+}), Q(1^{-})]$ follows from the limits of $Q$. For the
lower endpoint,
\begin{equation}
    Q_{\mathrm{RS}}(0^{+}) =
    \begin{cases}
        \lambda_1 - 1/\lambda_2, & \lambda_3 > 0, \\
        \lambda_1,               & \lambda_3 = 0, \\
        -\infty,                 & \lambda_3 < 0,
    \end{cases}
    \qquad
    Q_{\mathrm{FKML}}(0^{+}) =
    \begin{cases}
        \lambda_1 - 1/(\lambda_2\lambda_3), & \lambda_3 > 0, \\
        -\infty,                            & \lambda_3 \leq 0,
    \end{cases}
    \label{eq:support_lower}
\end{equation}
and for the upper endpoint,
\begin{equation}
    Q_{\mathrm{RS}}(1^{-}) =
    \begin{cases}
        \lambda_1 + 1/\lambda_2, & \lambda_4 > 0, \\
        \lambda_1,               & \lambda_4 = 0, \\
        +\infty,                 & \lambda_4 < 0,
    \end{cases}
    \qquad
    Q_{\mathrm{FKML}}(1^{-}) =
    \begin{cases}
        \lambda_1 + 1/(\lambda_2\lambda_4), & \lambda_4 > 0, \\
        +\infty,                            & \lambda_4 \leq 0,
    \end{cases}
    \label{eq:support_upper}
\end{equation}
where the infinite RS endpoints use $\lambda_2 < 0$, which the admissible
regions of Section~\ref{subsec:validity} require whenever a shape parameter is
negative. This is a fact worth keeping in view: \emph{a fitted GLD with
$\lambda_3, \lambda_4 > 0$ is bounded}. Approximating an unbounded target with
such a fit truncates its tails by construction, and the discrepancy that remains
is not a numerical artefact.

The first four moments exist when $\lambda_3, \lambda_4 > -1/4$, the pole of the
fourth moment. For the FKML form, they follow from the quantities
\begin{equation}
    v_k = \int_{0}^{1} [s(u)]^k \, du, \qquad k = 1, \ldots, 4,
\end{equation}
of the standardized quantile core
\begin{equation}
    s(u) = \frac{u^{\lambda_3}}{\lambda_3} - \frac{(1-u)^{\lambda_4}}{\lambda_4},
    \label{eq:su_definition}
\end{equation}
which are given explicitly by
\begin{equation}
    v_k = \sum_{j=0}^{k} \binom{k}{j} \frac{(-1)^j}{\lambda_3^{k-j} \lambda_4^j}
    B\big(\lambda_3(k-j)+1, \, \lambda_4 j+1\big),
    \label{eq:vk_beta}
\end{equation}
where $B(a, b)$ denotes the Euler beta function. From these, the mean, variance,
skewness and Pearson kurtosis are obtained in closed form; the RS form uses the
analogous coefficients $A$, $B$, $C$ and $D$ of \textcite{RambergSchmeiser1974},
with the skewness carrying the sign of $\lambda_2$. Both sets are implemented in
the package.

\subsection{Validity}\label{subsec:validity}

A parameter vector is a distribution only if $Q$ is non-decreasing on $(0,1)$.
Differentiating Equations~\eqref{eq:qrs} and \eqref{eq:qfkml},
\begin{equation}
    Q'_{\mathrm{RS}}(u) =
    \frac{\lambda_3 u^{\lambda_3-1} + \lambda_4 (1-u)^{\lambda_4-1}}{\lambda_2},
    \qquad
    Q'_{\mathrm{FKML}}(u) =
    \frac{u^{\lambda_3-1} + (1-u)^{\lambda_4-1}}{\lambda_2}.
    \label{eq:qprime}
\end{equation}
In the FKML form both terms of the numerator are positive for \emph{any} real
$\lambda_3, \lambda_4$, so $\lambda_2 > 0$ alone guarantees a valid
distribution.

In the RS form each shape parameter multiplies its own term, so a negative
$\lambda_3$ or $\lambda_4$ contributes a term of the opposite sign, and $Q'$ can
change sign on $(0,1)$. The admissible parameter vectors form the regions
catalogued by \textcite{KarianDudewicz2000}. Within the moment domain
$\lambda_3, \lambda_4 > -1/4$ four of them remain:
\begin{itemize}
    \item $\lambda_2 > 0$ and $\lambda_3, \lambda_4 \geq 0$, not both zero;
    \item $\lambda_2 < 0$ and $\lambda_3, \lambda_4 \leq 0$, not both zero;
    \item $\lambda_2 < 0$, $-1 < \lambda_3 < 0$, $\lambda_4 > 1$ and
    \begin{equation*}
        \frac{(1-\lambda_3)^{1-\lambda_3}\,(\lambda_4-1)^{\lambda_4-1}}
             {(\lambda_4-\lambda_3)^{\lambda_4-\lambda_3}}
        < -\frac{\lambda_3}{\lambda_4};
    \end{equation*}
    \item the mirror image of the previous region, with the roles of
    $\lambda_3$ and $\lambda_4$ exchanged.
\end{itemize}
The moment equations themselves do not encode these conditions: the
coefficients $A$ to $D$ are finite, and the equations can be solved, for vectors
that lie outside every region and do not define a distribution.
Section~\ref{sec:software} describes how the estimator excludes them.
